% To merge into the real deck: replace the next line with
% % !TeX root = raytracing/slides/main.tex
% raytracing_preamble.tex
% Common styling for Ray Tracing presentation

\documentclass[10pt]{beamer}
\usetheme{metropolis}
\usefonttheme{professionalfonts}

% --- Color Definitions ---
\definecolor{PrimaryColor}{RGB}{33,52,72}         % Deep navy
\definecolor{SecondaryColor}{RGB}{84,119,146}     % Desaturated blue-gray
\definecolor{AccentColor}{RGB}{100, 156, 165}       % Soft steel blue

\definecolor{BackgroundColor}{RGB}{245,247,250}   % Very light cool gray/blue-tinted white
\definecolor{TextColor}{RGB}{40,55,70}            % Darker, cool-toned charcoal (slightly less saturated)
\definecolor{LightGray}{RGB}{220,225,230}         % Soft, cool light gray
\definecolor{DarkGray}{RGB}{70,90,105}            % Cool-toned dark gray (coherent with Primary/Secondary)

\definecolor{RayColor}{RGB}{230,150,80}           % Muted warm orange (less saturated than pure orange)
\definecolor{ObjectColor}{RGB}{128,101,160}       % Muted violet (adjusted purple to match cool palette)
\definecolor{LightColor}{RGB}{235,200,100}        % Soft golden yellow (to complement cool blues)

% --- Theme Customization ---
\setbeamercolor{background canvas}{bg=BackgroundColor}
\setbeamercolor{normal text}{fg=TextColor}
\setbeamercolor{frametitle}{bg=PrimaryColor, fg=white}
\setbeamercolor{section in toc}{fg=PrimaryColor}
\setbeamercolor{block title}{bg=PrimaryColor!80, fg=white}
\setbeamercolor{block body}{bg=PrimaryColor!10}
\setbeamercolor{alerted text}{fg=AccentColor}
\setbeamercolor{itemize item}{fg=PrimaryColor}
\setbeamercolor{itemize subitem}{fg=SecondaryColor}
\setbeamerfont{frametitle}{size=\large,series=\bfseries}

% --- Packages ---
\usepackage[utf8]{inputenc}
\usepackage{url}
\usepackage{booktabs}
\usepackage{amsmath, amssymb}
\usepackage{fontawesome5}
\usepackage{pifont}
\usepackage[most]{tcolorbox}
\tcbuselibrary{skins}
\usepackage{colortbl}
\usepackage{array}
\usepackage{tikz}
\usepackage{graphicx}
\usetikzlibrary{shapes.callouts, positioning, arrows.meta, shapes.geometric, shadows, calc, patterns, 3d, backgrounds, shadings}
\usepackage{adjustbox}
\usepackage{ragged2e}
\usepackage{pgfplots}
\usepackage{graphicx}
\usepackage{caption}
\usepackage{minted}
\usepackage{subcaption}
\pgfplotsset{compat=1.18}

% --- Custom Commands ---
\newcommand{\cmark}{\textcolor{SecondaryColor}{\ding{51}}}
\newcommand{\xmark}{\textcolor{AccentColor}{\ding{55}}}
\newcommand{\highlight}[1]{\textcolor{PrimaryColor}{\textbf{#1}}}
\newcommand{\raycolor}[1]{\textcolor{RayColor}{\textbf{#1}}}
\newcommand{\objectcolor}[1]{\textcolor{ObjectColor}{\textbf{#1}}}

% --- TikZ styles for ray tracing visualizations ---
\tikzstyle{process} = [rectangle, rounded corners=3mm, minimum width=2cm, minimum height=0.8cm, text centered, draw=PrimaryColor, thick, fill=PrimaryColor!15, drop shadow]
\tikzstyle{arrow} = [thick, PrimaryColor, ->, >=stealth]
\tikzstyle{ray} = [thick, RayColor, ->, >=stealth]
\tikzstyle{lightray} = [thick, LightColor, ->, >=stealth]
\tikzstyle{reflectray} = [thick, SecondaryColor, ->, >=stealth]
\tikzstyle{refractray} = [thick, AccentColor, ->, >=stealth]
\tikzstyle{shadowray} = [thick, DarkGray, ->, >=stealth, dashed]

% --- 3D object styles ---
\tikzstyle{sphere} = [circle, minimum size=1.5cm, draw=ObjectColor, thick, fill=ObjectColor!20, drop shadow]
\tikzstyle{plane} = [rectangle, minimum width=3cm, minimum height=0.2cm, draw=ObjectColor, thick, fill=ObjectColor!20]
\tikzstyle{triangle} = [regular polygon, regular polygon sides=3, minimum size=1.5cm, draw=ObjectColor, thick, fill=ObjectColor!20]

% --- Eye/Camera styles ---
\tikzstyle{eye} = [circle, minimum size=0.8cm, draw=PrimaryColor, thick, fill=PrimaryColor!30]
\tikzstyle{pixel} = [rectangle, minimum size=0.2cm, draw=AccentColor, fill=AccentColor!30]

% --- Text box styles ---
\tikzstyle{conceptbox} = [rectangle, rounded corners, fill=PrimaryColor!10, draw=PrimaryColor, thick, text width=0.8\textwidth, inner sep=8pt]
\tikzstyle{formulabox} = [rectangle, rounded corners, fill=SecondaryColor!10, draw=SecondaryColor, thick, inner sep=6pt]

% --- Custom environments ---
\newtcolorbox{raybox}[1]{
  colback=RayColor!10,
  colframe=RayColor,
  title=#1,
  fonttitle=\bfseries,
  sharp corners
}

\newtcolorbox{conceptbox}[1]{
  colback=PrimaryColor!10,
  colframe=PrimaryColor,
  title=#1,
  fonttitle=\bfseries,
  rounded corners
}

\newtcolorbox{mathbox}[1]{
  colback=SecondaryColor!10,
  colframe=SecondaryColor,
  title=#1,
  fonttitle=\bfseries,
  rounded corners
}

\newenvironment{timeline}{%
  \begin{tikzpicture}[scale=0.8]
    \coordinate (start) at (0,0);
    \newcounter{timelineitem}
    \setcounter{timelineitem}{0}
  }{%
  \end{tikzpicture}
}

\newcommand{\timelineitem}[4]{%
  \stepcounter{timelineitem}%
  % Name of this item’s node
  \edef\thisID{time\thetimelineitem}%
  \ifnum\value{timelineitem}=1
  % First item: anchor at (start)
  \node[rectangle, draw, fill=PrimaryColor!20,
  text width=1.5cm, minimum height=0.8cm, align=center]
  (\thisID) at (start) {\textbf{#2}};
  \else
  % Later items: below=#1 of previous time node
  \pgfmathtruncatemacro\prev{\value{timelineitem}-1}%
  \node[rectangle, draw, fill=PrimaryColor!20,
    text width=1.5cm, minimum height=0.8cm, align=center,
  below=#1 of time\prev]
  (\thisID) {\textbf{#2}};
  \fi
  % Description always to the right of this time node
  \node[rectangle, draw, fill=SecondaryColor!10,
    text width=4cm, minimum height=0.8cm,
  align=left, right=0.3cm of \thisID]
  (desc\thetimelineitem) {\textbf{#3}\\#4};
  % Draw arrow back to previous if not the first
  \ifnum\value{timelineitem}>1
  \draw[->, thick, PrimaryColor]
  (time\prev.south) -- (\thisID.north);
  \fi
}

\tikzset{
  camera/.style={fill=PrimaryColor!60, draw=PrimaryColor!80, rectangle, minimum size=8pt},
  image plane/.style={fill=AccentColor!10, draw=AccentColor!50, opacity=0.8},
  pixel/.style={fill=AccentColor!60, thick},
  primary ray/.style={->, very thick, red!90},
  object/.style={fill=ObjectColor!60, draw=ObjectColor!80, circle, minimum size=12pt},
  fovangle/.style={<->, thick, PrimaryColor, dashed}
}

\tikzset{
  lens/.style={thick, PrimaryColor, line width=3pt},
  focal plane/.style={thick, AccentColor},
  object ray/.style={->, thick, ObjectColor},
  image ray/.style={->, thick, SecondaryColor},
  optical axis/.style={dashed, gray}
}

\newcommand{\lens}[3]{%
  \begingroup
  % half‐dimensions
  \pgfmathsetlengthmacro{\a}{#2}%
  \pgfmathsetlengthmacro{\b}{#3}%
  \pgfmathsetlengthmacro{\c}{#3/2}%
  \begin{scope}[shift={#1}]
    \draw[line join=round, fill=blue!15]
    (0,-{\c})
    arc(-30:30:{\a} and {\b})
    arc(150:210:{\a} and {\b})
    ;
  \end{scope}
  \endgroup
}

% ================================================================
% Reference theme, rebuilt to match combined/deck.pdf
% Rust title bar, cream ground, Adventor headings, thin footer.
% Placed last so it overrides metropolis.
% ================================================================
\usepackage{tgheros}
\usepackage{tgadventor}
\renewcommand{\sfdefault}{qhv}                       % Heros for body copy
\newcommand{\headfont}{\fontfamily{qag}\selectfont}  % Adventor for headings

\definecolor{ThemeRust}{HTML}{9B4632}
\definecolor{ThemeCream}{HTML}{F8F6F2}
\definecolor{ThemeSage}{HTML}{C6CCB1}
\definecolor{ThemeStone}{HTML}{E7E7DB}
\definecolor{ThemeMuted}{HTML}{8B8A88}
\definecolor{ThemeInk}{HTML}{191919}

\setbeamercolor{background canvas}{bg=ThemeCream}
\setbeamercolor{normal text}{fg=ThemeInk}
\setbeamercolor{frametitle}{bg=ThemeRust, fg=ThemeCream}
\setbeamercolor{footer}{bg=ThemeCream, fg=ThemeMuted}
\setbeamercolor{alerted text}{fg=ThemeRust}
\setbeamercolor{itemize item}{fg=ThemeRust}
\setbeamercolor{itemize subitem}{fg=ThemeRust}
\setbeamercolor{palette primary}{bg=ThemeRust, fg=ThemeCream}  % section pages

% full-bleed rust bar across the top
\setbeamertemplate{frametitle}{%
  \nointerlineskip%
  \begin{beamercolorbox}[wd=\paperwidth,sep=0pt]{frametitle}%
    \vbox{%
      \vskip0.60em
      \hbox{\hskip0.9em%
        \begin{minipage}{\dimexpr\paperwidth-1.8em\relax}%
          \raggedright\headfont\bfseries\large\insertframetitle%
        \end{minipage}}%
      \vskip0.55em
    }%
  \end{beamercolorbox}%
}

% thin three-zone footer
\setbeamertemplate{footline}{%
  \hbox{%
  \begin{beamercolorbox}[wd=.40\paperwidth,ht=1.9ex,dp=1.0ex,left]{footer}%
    \hspace*{0.9em}\tiny\insertshortauthor
  \end{beamercolorbox}%
  \begin{beamercolorbox}[wd=.28\paperwidth,ht=1.9ex,dp=1.0ex,center]{footer}%
    \tiny\insertshorttitle
  \end{beamercolorbox}%
  \begin{beamercolorbox}[wd=.32\paperwidth,ht=1.9ex,dp=1.0ex,right]{footer}%
    \tiny\insertframenumber\,/\,\inserttotalframenumber\hspace*{0.9em}
  \end{beamercolorbox}}%
  \vskip1pt
}

% full-bleed rust title page
\setbeamertemplate{title page}{%
  \begin{tikzpicture}[remember picture, overlay]
    \fill[ThemeRust] (current page.south west) rectangle (current page.north east);
  \end{tikzpicture}%
  \vfill
  \begin{flushright}
    {\headfont\bfseries\LARGE\color{ThemeCream}\inserttitle\par}
    \vspace{0.3em}
    {\color{ThemeCream!85}\small\insertsubtitle\par}
    \vspace{1.3em}
    {\color{ThemeCream}\small\insertauthor\par}
    \vspace{0.5em}
    {\color{ThemeCream!80}\scriptsize\insertinstitute\par}
    \vspace{0.8em}
    {\color{ThemeCream!75}\scriptsize\insertdate\par}
  \end{flushright}
  \vfill
}

% --- Deck palette remapped onto the theme -------------------------------
% Overrides the navy/violet set defined at the top of this file. Every TikZ
% figure and tcolorbox references these names, so they all follow along.
\definecolor{PrimaryColor}{HTML}{9B4632}     % rust: highlights, frames, titles
\definecolor{SecondaryColor}{HTML}{7A5C4B}   % warm brown: secondary boxes
\definecolor{AccentColor}{HTML}{7E8C57}      % sage: arrows, alerted text
\definecolor{ObjectColor}{HTML}{A9745B}      % clay: solid shapes
\definecolor{LightColor}{HTML}{D9A441}       % warm gold
\definecolor{TextColor}{HTML}{191919}
\definecolor{BackgroundColor}{HTML}{F8F6F2}
\definecolor{LightGray}{HTML}{E7E7DB}
\definecolor{DarkGray}{HTML}{6B6257}

% the four fern rules -- must match COLS in part 3/figs5.py
\definecolor{RuleA}{HTML}{C0392B}
\definecolor{RuleB}{HTML}{2E7D32}
\definecolor{RuleC}{HTML}{B8860B}
\definecolor{RuleD}{HTML}{1F6F8B}

  (the teacher's actual preamble)
% !TeX root = raytracing/slides/main.tex
% raytracing_preamble.tex
% Common styling for Ray Tracing presentation

\documentclass[10pt]{beamer}
\usetheme{metropolis}
\usefonttheme{professionalfonts}

% --- Color Definitions ---
\definecolor{PrimaryColor}{RGB}{33,52,72}         % Deep navy
\definecolor{SecondaryColor}{RGB}{84,119,146}     % Desaturated blue-gray
\definecolor{AccentColor}{RGB}{100, 156, 165}       % Soft steel blue

\definecolor{BackgroundColor}{RGB}{245,247,250}   % Very light cool gray/blue-tinted white
\definecolor{TextColor}{RGB}{40,55,70}            % Darker, cool-toned charcoal (slightly less saturated)
\definecolor{LightGray}{RGB}{220,225,230}         % Soft, cool light gray
\definecolor{DarkGray}{RGB}{70,90,105}            % Cool-toned dark gray (coherent with Primary/Secondary)

\definecolor{RayColor}{RGB}{230,150,80}           % Muted warm orange (less saturated than pure orange)
\definecolor{ObjectColor}{RGB}{128,101,160}       % Muted violet (adjusted purple to match cool palette)
\definecolor{LightColor}{RGB}{235,200,100}        % Soft golden yellow (to complement cool blues)

% --- Theme Customization ---
\setbeamercolor{background canvas}{bg=BackgroundColor}
\setbeamercolor{normal text}{fg=TextColor}
\setbeamercolor{frametitle}{bg=PrimaryColor, fg=white}
\setbeamercolor{section in toc}{fg=PrimaryColor}
\setbeamercolor{block title}{bg=PrimaryColor!80, fg=white}
\setbeamercolor{block body}{bg=PrimaryColor!10}
\setbeamercolor{alerted text}{fg=AccentColor}
\setbeamercolor{itemize item}{fg=PrimaryColor}
\setbeamercolor{itemize subitem}{fg=SecondaryColor}
\setbeamerfont{frametitle}{size=\large,series=\bfseries}

% --- Packages ---
\usepackage[utf8]{inputenc}
\usepackage{url}
\usepackage{booktabs}
\usepackage{amsmath, amssymb}
\usepackage{fontawesome5}
\usepackage{pifont}
\usepackage[most]{tcolorbox}
\tcbuselibrary{skins}
\usepackage{colortbl}
\usepackage{array}
\usepackage{tikz}
\usepackage{graphicx}
\usetikzlibrary{shapes.callouts, positioning, arrows.meta, shapes.geometric, shadows, calc, patterns, 3d, backgrounds, shadings}
\usepackage{adjustbox}
\usepackage{ragged2e}
\usepackage{pgfplots}
\usepackage{graphicx}
\usepackage{caption}
\usepackage{minted}
\usepackage{subcaption}
\pgfplotsset{compat=1.18}

% --- Custom Commands ---
\newcommand{\cmark}{\textcolor{SecondaryColor}{\ding{51}}}
\newcommand{\xmark}{\textcolor{AccentColor}{\ding{55}}}
\newcommand{\highlight}[1]{\textcolor{PrimaryColor}{\textbf{#1}}}
\newcommand{\raycolor}[1]{\textcolor{RayColor}{\textbf{#1}}}
\newcommand{\objectcolor}[1]{\textcolor{ObjectColor}{\textbf{#1}}}

% --- TikZ styles for ray tracing visualizations ---
\tikzstyle{process} = [rectangle, rounded corners=3mm, minimum width=2cm, minimum height=0.8cm, text centered, draw=PrimaryColor, thick, fill=PrimaryColor!15, drop shadow]
\tikzstyle{arrow} = [thick, PrimaryColor, ->, >=stealth]
\tikzstyle{ray} = [thick, RayColor, ->, >=stealth]
\tikzstyle{lightray} = [thick, LightColor, ->, >=stealth]
\tikzstyle{reflectray} = [thick, SecondaryColor, ->, >=stealth]
\tikzstyle{refractray} = [thick, AccentColor, ->, >=stealth]
\tikzstyle{shadowray} = [thick, DarkGray, ->, >=stealth, dashed]

% --- 3D object styles ---
\tikzstyle{sphere} = [circle, minimum size=1.5cm, draw=ObjectColor, thick, fill=ObjectColor!20, drop shadow]
\tikzstyle{plane} = [rectangle, minimum width=3cm, minimum height=0.2cm, draw=ObjectColor, thick, fill=ObjectColor!20]
\tikzstyle{triangle} = [regular polygon, regular polygon sides=3, minimum size=1.5cm, draw=ObjectColor, thick, fill=ObjectColor!20]

% --- Eye/Camera styles ---
\tikzstyle{eye} = [circle, minimum size=0.8cm, draw=PrimaryColor, thick, fill=PrimaryColor!30]
\tikzstyle{pixel} = [rectangle, minimum size=0.2cm, draw=AccentColor, fill=AccentColor!30]

% --- Text box styles ---
\tikzstyle{conceptbox} = [rectangle, rounded corners, fill=PrimaryColor!10, draw=PrimaryColor, thick, text width=0.8\textwidth, inner sep=8pt]
\tikzstyle{formulabox} = [rectangle, rounded corners, fill=SecondaryColor!10, draw=SecondaryColor, thick, inner sep=6pt]

% --- Custom environments ---
\newtcolorbox{raybox}[1]{
  colback=RayColor!10,
  colframe=RayColor,
  title=#1,
  fonttitle=\bfseries,
  sharp corners
}

\newtcolorbox{conceptbox}[1]{
  colback=PrimaryColor!10,
  colframe=PrimaryColor,
  title=#1,
  fonttitle=\bfseries,
  rounded corners
}

\newtcolorbox{mathbox}[1]{
  colback=SecondaryColor!10,
  colframe=SecondaryColor,
  title=#1,
  fonttitle=\bfseries,
  rounded corners
}

\newenvironment{timeline}{%
  \begin{tikzpicture}[scale=0.8]
    \coordinate (start) at (0,0);
    \newcounter{timelineitem}
    \setcounter{timelineitem}{0}
  }{%
  \end{tikzpicture}
}

\newcommand{\timelineitem}[4]{%
  \stepcounter{timelineitem}%
  % Name of this item’s node
  \edef\thisID{time\thetimelineitem}%
  \ifnum\value{timelineitem}=1
  % First item: anchor at (start)
  \node[rectangle, draw, fill=PrimaryColor!20,
  text width=1.5cm, minimum height=0.8cm, align=center]
  (\thisID) at (start) {\textbf{#2}};
  \else
  % Later items: below=#1 of previous time node
  \pgfmathtruncatemacro\prev{\value{timelineitem}-1}%
  \node[rectangle, draw, fill=PrimaryColor!20,
    text width=1.5cm, minimum height=0.8cm, align=center,
  below=#1 of time\prev]
  (\thisID) {\textbf{#2}};
  \fi
  % Description always to the right of this time node
  \node[rectangle, draw, fill=SecondaryColor!10,
    text width=4cm, minimum height=0.8cm,
  align=left, right=0.3cm of \thisID]
  (desc\thetimelineitem) {\textbf{#3}\\#4};
  % Draw arrow back to previous if not the first
  \ifnum\value{timelineitem}>1
  \draw[->, thick, PrimaryColor]
  (time\prev.south) -- (\thisID.north);
  \fi
}

\tikzset{
  camera/.style={fill=PrimaryColor!60, draw=PrimaryColor!80, rectangle, minimum size=8pt},
  image plane/.style={fill=AccentColor!10, draw=AccentColor!50, opacity=0.8},
  pixel/.style={fill=AccentColor!60, thick},
  primary ray/.style={->, very thick, red!90},
  object/.style={fill=ObjectColor!60, draw=ObjectColor!80, circle, minimum size=12pt},
  fovangle/.style={<->, thick, PrimaryColor, dashed}
}

\tikzset{
  lens/.style={thick, PrimaryColor, line width=3pt},
  focal plane/.style={thick, AccentColor},
  object ray/.style={->, thick, ObjectColor},
  image ray/.style={->, thick, SecondaryColor},
  optical axis/.style={dashed, gray}
}

\newcommand{\lens}[3]{%
  \begingroup
  % half‐dimensions
  \pgfmathsetlengthmacro{\a}{#2}%
  \pgfmathsetlengthmacro{\b}{#3}%
  \pgfmathsetlengthmacro{\c}{#3/2}%
  \begin{scope}[shift={#1}]
    \draw[line join=round, fill=blue!15]
    (0,-{\c})
    arc(-30:30:{\a} and {\b})
    arc(150:210:{\a} and {\b})
    ;
  \end{scope}
  \endgroup
}

% ================================================================
% Reference theme, rebuilt to match combined/deck.pdf
% Rust title bar, cream ground, Adventor headings, thin footer.
% Placed last so it overrides metropolis.
% ================================================================
\usepackage{tgheros}
\usepackage{tgadventor}
\renewcommand{\sfdefault}{qhv}                       % Heros for body copy
\newcommand{\headfont}{\fontfamily{qag}\selectfont}  % Adventor for headings

\definecolor{ThemeRust}{HTML}{9B4632}
\definecolor{ThemeCream}{HTML}{F8F6F2}
\definecolor{ThemeSage}{HTML}{C6CCB1}
\definecolor{ThemeStone}{HTML}{E7E7DB}
\definecolor{ThemeMuted}{HTML}{8B8A88}
\definecolor{ThemeInk}{HTML}{191919}

\setbeamercolor{background canvas}{bg=ThemeCream}
\setbeamercolor{normal text}{fg=ThemeInk}
\setbeamercolor{frametitle}{bg=ThemeRust, fg=ThemeCream}
\setbeamercolor{footer}{bg=ThemeCream, fg=ThemeMuted}
\setbeamercolor{alerted text}{fg=ThemeRust}
\setbeamercolor{itemize item}{fg=ThemeRust}
\setbeamercolor{itemize subitem}{fg=ThemeRust}
\setbeamercolor{palette primary}{bg=ThemeRust, fg=ThemeCream}  % section pages

% full-bleed rust bar across the top
\setbeamertemplate{frametitle}{%
  \nointerlineskip%
  \begin{beamercolorbox}[wd=\paperwidth,sep=0pt]{frametitle}%
    \vbox{%
      \vskip0.60em
      \hbox{\hskip0.9em%
        \begin{minipage}{\dimexpr\paperwidth-1.8em\relax}%
          \raggedright\headfont\bfseries\large\insertframetitle%
        \end{minipage}}%
      \vskip0.55em
    }%
  \end{beamercolorbox}%
}

% thin three-zone footer
\setbeamertemplate{footline}{%
  \hbox{%
  \begin{beamercolorbox}[wd=.40\paperwidth,ht=1.9ex,dp=1.0ex,left]{footer}%
    \hspace*{0.9em}\tiny\insertshortauthor
  \end{beamercolorbox}%
  \begin{beamercolorbox}[wd=.28\paperwidth,ht=1.9ex,dp=1.0ex,center]{footer}%
    \tiny\insertshorttitle
  \end{beamercolorbox}%
  \begin{beamercolorbox}[wd=.32\paperwidth,ht=1.9ex,dp=1.0ex,right]{footer}%
    \tiny\insertframenumber\,/\,\inserttotalframenumber\hspace*{0.9em}
  \end{beamercolorbox}}%
  \vskip1pt
}

% full-bleed rust title page
\setbeamertemplate{title page}{%
  \begin{tikzpicture}[remember picture, overlay]
    \fill[ThemeRust] (current page.south west) rectangle (current page.north east);
  \end{tikzpicture}%
  \vfill
  \begin{flushright}
    {\headfont\bfseries\LARGE\color{ThemeCream}\inserttitle\par}
    \vspace{0.3em}
    {\color{ThemeCream!85}\small\insertsubtitle\par}
    \vspace{1.3em}
    {\color{ThemeCream}\small\insertauthor\par}
    \vspace{0.5em}
    {\color{ThemeCream!80}\scriptsize\insertinstitute\par}
    \vspace{0.8em}
    {\color{ThemeCream!75}\scriptsize\insertdate\par}
  \end{flushright}
  \vfill
}

% --- Deck palette remapped onto the theme -------------------------------
% Overrides the navy/violet set defined at the top of this file. Every TikZ
% figure and tcolorbox references these names, so they all follow along.
\definecolor{PrimaryColor}{HTML}{9B4632}     % rust: highlights, frames, titles
\definecolor{SecondaryColor}{HTML}{7A5C4B}   % warm brown: secondary boxes
\definecolor{AccentColor}{HTML}{7E8C57}      % sage: arrows, alerted text
\definecolor{ObjectColor}{HTML}{A9745B}      % clay: solid shapes
\definecolor{LightColor}{HTML}{D9A441}       % warm gold
\definecolor{TextColor}{HTML}{191919}
\definecolor{BackgroundColor}{HTML}{F8F6F2}
\definecolor{LightGray}{HTML}{E7E7DB}
\definecolor{DarkGray}{HTML}{6B6257}

% the four fern rules -- must match COLS in part 3/figs5.py
\definecolor{RuleA}{HTML}{C0392B}
\definecolor{RuleB}{HTML}{2E7D32}
\definecolor{RuleC}{HTML}{B8860B}
\definecolor{RuleD}{HTML}{1F6F8B}



\title{Fractals}
\subtitle{Infinite Details from Simple Rules}
% ===== SPEAKER NAMES - the only place to edit them =====
% \author[short form -> page footer]{full form -> title page}
\author[Aditya, Sami, Pritom]{%
  Aditya Tirtho Roy (2305178)\\
  Farhan Ehsas Sami (2305177)\\
  Pritom Biswas (2305160)}
\institute{Department of Computer Science and Engineering\\
  Bangladesh University of Engineering and Technology (BUET)}
\date{CSE 200 Presentation}

% one-line takeaway box used on every slide
\newcommand{\takeaway}[1]{%
  \begin{tcolorbox}[colback=PrimaryColor!8, colframe=PrimaryColor, arc=2mm,
    boxrule=0.8pt, left=6pt, right=6pt, top=3pt, bottom=3pt]
    \small #1
  \end{tcolorbox}}
% small muted source line
\newcommand{\source}[1]{{\scriptsize\color{DarkGray}#1}}


\begin{document}

\begin{frame}[plain]
  \titlepage
\end{frame}

% ================================================================
% PART 1 — HISTORY (Aditya): the hook, the triangle, the
% coastline puzzle, then the compressed flashback.
% ================================================================
\section{History}


% ---------------------------------------------------------------
% Frame 1 of 4: The hook — designing trees the hard way vs. one rule
% ---------------------------------------------------------------
\begin{frame}{So, You're Designing All the Trees\ldots}
  \small
  \onslide<1->{Imagine you're a game developer. You've just been
  handed one job: design \highlight{every tree} in the game ---
  all of them.}

  \vspace{0.2cm}
  \begin{columns}[c]
    \begin{column}{0.55\textwidth}
      \onslide<2->{You could hand-draw every branch, every twig,
      every leaf, one at a time\ldots}

      \vspace{0.15cm}
      \onslide<3->{\itshape (kind of like how Rockstar's been
      ``finishing'' GTA~6 for the last decade.)}
    \end{column}
    \begin{column}{0.4\textwidth}
      \centering
      \onslide<4->{\includegraphics[height=2.2cm]{images/tree/img17-removebg-preview.png}}
    \end{column}
  \end{columns}

  \onslide<4->{
    \vspace{0.12cm}
    \begin{center}
      \begin{tikzpicture}
        \node[conceptbox, text width=0.85\textwidth] {
          \centering
          Or, you write \highlight{one rule} --- ``a branch splits
          into two smaller branches'' --- and let that rule call
          itself, again and again, at a smaller scale every time.
        };
      \end{tikzpicture}
    \end{center}
  }

  \onslide<5->{
    \vspace{0.06cm}
    \takeaway{Trunk becomes branch, branch becomes twig, twig
    becomes leaf --- \highlight{same rule, every time}. The engine
    genuinely can't tell the whole tree apart from one tiny piece
    of it.}
  }
\end{frame}

% ---------------------------------------------------------------
% Frame 2 of 4: The Sierpinski teaser — same trick, simplest form
% ---------------------------------------------------------------
\begin{frame}{D\'ej\`a Vu?}
  \begin{columns}[c]
    \begin{column}{0.55\textwidth}
      \small
      \onslide<1->{Here's the simplest possible version of that
      same trick: draw a triangle.}

      \vspace{0.15cm}
      \onslide<2->{Cut a smaller triangle out of the middle.}

      \vspace{0.15cm}
      \onslide<3->{Now do that again, inside each of the three
      triangles you've got left.}

      \vspace{0.15cm}
      \onslide<4->{And again --- forever.}

      \vspace{0.3cm}
      \onslide<5->{
        \begin{tikzpicture}
          \node[conceptbox, text width=0.9\textwidth] {
            \centering
            It's basically filming yourself watching TV, while the
            TV plays that exact clip of you watching TV.
          };
        \end{tikzpicture}
      }
    \end{column}
    \begin{column}{0.4\textwidth}
      \centering
      \foreach \t/\i in {1/0,2/1,3/2,4/3} {
        \only<\t>{
          \includegraphics[width=0.85\textwidth]{images/sierpinski/\i.png}
        }
      }
      \only<5->{\includegraphics[width=0.85\textwidth]{images/sierpinski/tvwatch.jpg}}
    \end{column}
  \end{columns}

  \onslide<6->{
    \vspace{0.2cm}
    \takeaway{Weird, right? That's a \highlight{fractal} --- a
    shape built out of smaller copies of itself.}
  }
\end{frame}

% ---------------------------------------------------------------
% Moved up from the Definition section: the Deja Vu frame ends on
% "a shape built out of smaller copies of itself", so the formal
% definition belongs here rather than after the history.
% ---------------------------------------------------------------
\begin{frame}{What are Fractals?}
  \begin{conceptbox}{Definition}
    A \highlight{fractal} is a geometric shape containing detailed
    structure at arbitrarily small scales, usually having a
    \highlight{fractal dimension} strictly exceeding the
    \highlight{topological dimension}.
  \end{conceptbox}

  \vspace{0.1cm}

  \textbf{Key Characteristics:}
  \begin{itemize}
    \item \highlight{Scale invariance}: Details at every level of magnification
    \item \highlight{Fractal dimension}: More than the \textit{conventional} dimension
    \item \highlight{Infinite complexity}: Generated by simple, recursive rules
      \begin{itemize}
        \item Like the infinite complexities of nature
      \end{itemize}
  \end{itemize}

  \vspace{0.1cm}

  \pause
  \begin{center}
    \begin{tikzpicture}
      \node[conceptbox, text width=0.7\textwidth] {
        \centering
        \textit{"Clouds are not spheres, mountains are not cones, coastlines are not circles..."} \\
        \textbf{--- Benoit Mandelbrot}
      };
    \end{tikzpicture}
  \end{center}
\end{frame}

% ---------------------------------------------------------------
% Frame 3 of 4: Mandelbrot's question — the coastline pulls history
% in on its own, instead of a separate "meanwhile, in the 1800s..."
% ---------------------------------------------------------------
\begin{frame}{Mandelbrot's Question}
  \small
  In the 1960s, working at IBM, \highlight{Benoit Mandelbrot} asked a
  deceptively simple question:

  \begin{center}
    \begin{tikzpicture}
      \node[conceptbox, text width=0.8\textwidth] {
        \centering
        \textit{``How long is the coast of Britain?''}
      };
    \end{tikzpicture}
  \end{center}

  \pause
  \vspace{0.2cm}
  The honest answer: \highlight{it depends on your ruler}. Measure
  with a 100~km ruler, get one length. Use a 1~km ruler, and the
  measured length grows — every bay and inlet you'd smoothed over
  now counts. The length never converges to a single number.

  \pause
  \vspace{0.2cm}
  Mandelbrot's insight: this isn't a flaw in the coastline — it's a
  sign that classical length and area are \highlight{the wrong tools}
  for something like this. What a coastline actually needed was
  exactly the kind of \highlight{``useless'' geometry} mathematicians
  had quietly been building for a century already --- without knowing
  why.
\end{frame}

% ---------------------------------------------------------------
% Frame 4 of 4: Where that toolkit came from — compressed flashback,
% deliberately ends back on the Sierpinski triangle for the handoff.
% ---------------------------------------------------------------
\begin{frame}{Borrowed Tools}
  \small
  Mandelbrot wasn't building from nothing.

  \pause
  \vspace{0.2cm}
  \begin{center}
    \begin{tikzpicture}
      \node[conceptbox, text width=0.8\textwidth] {
        \centering
        Back in 1872, \highlight{Karl Weierstrass} broke an
        assumption every mathematician had taken for granted ---
        on purpose.
      };
    \end{tikzpicture}
  \end{center}

  \pause
  \vspace{0.2cm}
  That set off a decades-long spree: mathematicians started building
  weirder and weirder shapes on purpose, just to see where ordinary
  geometry broke down --- \highlight{Cantor}'s dust and
  \highlight{Koch}'s endless jagged curve.

  \pause
  \vspace{0.2cm}
  For nearly a century, these were dismissed as
  \highlight{beautiful but useless}. Mandelbrot was the one who
  realized they weren't useless at all --- he gave the whole family
  a name: \highlight{fractal}.

  \pause
  \vspace{0.3cm}
  \begin{center}
    \begin{tikzpicture}
      \node[conceptbox, text width=0.85\textwidth] {
        \centering
        And that triangle? The exact one you watched me draw a few
        minutes ago.
      };
    \end{tikzpicture}
  \end{center}
\end{frame}

% ================================================================
% PART 2 — SIERPINSKI CONSTRUCTION WALKTHROUGH
% ================================================================
\section{Some Fractals}

\begin{frame}{The Sierpinski Triangle}
  \begin{columns}
    \begin{column}{0.6\textwidth}
      \small
      \only<1-4>{
        Let's construct the Sierpinski Triangle:
        \begin{conceptbox}{Construction Process}
          \footnotesize
          \begin{enumerate}
            \item Start with an equilateral triangle
            \item<2-> Remove the central triangle (triangle connecting midpoints of the sides)
            \item<3-> Repeat for each triangle
            \item<4-> After infinite iterations, we get the Sierpinski Triangle
          \end{enumerate}
        \end{conceptbox}
      }
    \end{column}
    \begin{column}{0.4\textwidth}
      \begin{center}
        \begin{figure}
          \foreach \t/\i in {1/0,2/1,3/2,4/3} {
            \only<\t>{
              \includegraphics[width=\textwidth]{images/sierpinski/\i.png}
              \caption*{Iteration \i}
            }
          }
        \end{figure}
      \end{center}
    \end{column}
  \end{columns}
\end{frame}

% ================================================================
% PART 3 — SIERPINSKI AREA & PERIMETER CALCULATION
% ================================================================
\section{Area \& Perimeter}

\begin{frame}{The Sierpinski Triangle: Area}
  \begin{columns}
    \begin{column}{0.6\textwidth}
      \small
      Now let's find the area of triangle(s):
      \begin{mathbox}{Area Calculation}
        \footnotesize
        \begin{align*}
          \only<1->{
            A_0 &= \frac{1}{2}\cdot \sin(60^\circ) \cdot a^2 = \frac{\sqrt{3}}{4} a^2 \\
          }
          \only<2>{
            A_1 &= A_0 - \frac{1}{2}\cdot \sin(60^\circ) \cdot \left(\frac{a}{2}\right)^2 \\
            &= A_0 - \frac{\sqrt{3}}{16} a^2 = A_0 - \frac{1}{4} A_0 \\
            &= \frac{3}{4} A_0
          }
          \only<3->{
            A_1  &= \left(1-\frac{1}{4}\right) A_0 =  \frac{3}{4} A_0 \\
            A_2 &= \left(1-\frac{1}{4}\right) A_1 = \left(\frac{3}{4}\right)^2 A_0 \\
          }
          \only<4->{
            \hdots \\
            A_n &= \left(1-\frac{1}{4}\right)^n A_0 = \left(\frac{3}{4}\right)^n A_0 \\
            \lim_{n \to \infty} A_n &= 0
          }
        \end{align*}
      \end{mathbox}
      \only<2-3>{
        We cut out one fourth of the area at each step. So, naturally
        area is reduced by $\frac{3}{4}$ at each step.
      }
      \only<4->{
        \highlight{The area of the Sierpinski Triangle is zero even
        though it is a 2D shape.}
      }
    \end{column}
    \begin{column}{0.4\textwidth}
      \begin{center}
        \begin{figure}
          \foreach \t/\i in {1/0, 2/1, 3/2, 4/3} {
            \only<\t>{
              \includegraphics[width=\textwidth]{images/sierpinski/\i.png}
              \caption*{Iteration \i}
            }
          }
        \end{figure}
      \end{center}
    \end{column}
  \end{columns}
\end{frame}

\begin{frame}{The Sierpinski Triangle: Perimeter}
  \begin{columns}
    \begin{column}{0.6\textwidth}
      \small
      Now let's find the perimeter:
      \begin{mathbox}{Perimeter Calculation}
        \footnotesize
        \begin{align*}
          \only<1->{
            P_0 &= 3a \\
          }
          \only<2->{
            P_1 &= P_0 + 3\cdot\frac{a}{2} = \frac{9}{2} a = \frac{3}{2} P_0 \\
          }
          \only<3->{
            P_2 &= P_1 + 3\cdot3\cdot\frac{a}{4} = P_1 + \frac{1}{2} P_1 = \frac{3}{2} P_1 \\
          }
          \only<4->{
            \hdots \\
            P_n &= \frac{3}{2} P_{n-1} = \left(\frac{3}{2}\right)^n P_0 \\
            \lim_{n \to \infty} P_n &= \infty
          }
        \end{align*}
      \end{mathbox}
      \only<2-3>{
        We add half of the previous perimeter at each step.
        This is because we add the same number of triangles as the
        last step, but each triangle is half the size of the
        previous one.

        As a result, the perimeter grows by a factor of $\frac{3}{2}$
        at each step.
      }
      \only<4->{
        \highlight{The Sierpinski Triangle has infinite perimeter but
        zero area.}
      }
    \end{column}
    \begin{column}{0.4\textwidth}
      \begin{center}
        \begin{figure}
          \foreach \t/\i in {1/0, 2/1, 3/2, 4/3} {
            \only<\t>{
              \includegraphics[width=\textwidth]{images/sierpinski/\i.png}
              \caption*{Iteration \i}
            }
          }
        \end{figure}
      \end{center}
    \end{column}
  \end{columns}
\end{frame}

% ================================================================
% Handoff from Speaker 1
% ================================================================
\begin{frame}{Handoff to Speaker 2}
  \small
  We've now seen something strange, and precise:

  \begin{itemize}
    \item A \highlight{2D shape} (Sierpinski Triangle) with
      \highlight{zero area}.
    \item A shape built from a rule so simple a child could apply it
      by hand.
    \item No matter how far we zoom in, there's always more detail.
  \end{itemize}

  \pause
  \vspace{0.3cm}
  \begin{center}
    \begin{tikzpicture}
      \node[conceptbox, text width=0.8\textwidth] {
        \centering
        If it isn't a line, and it isn't a surface --- what
        \highlight{is} its true dimension?
      };
    \end{tikzpicture}
  \end{center}

  \vspace{0.2cm}
  \pause
  That number is exactly what Hausdorff was reaching for in 1918 ---
  and it's where Speaker~2 picks up the story.
\end{frame}

% ================================================================
% PART 2 — FRACTAL DIMENSION (Speaker 2)
% 9 frames: transition, integer dims, Sierpinski reveal,
%   scaling, formula, Sierpinski derivation, Koch derivation,
%   box-counting theory, handoff to Speaker 3
% ================================================================
\section{Fractal Dimension}

% ---------------------------------------------------------------
% S2 Frame 1: Transition — section title
% ---------------------------------------------------------------
\begin{frame}{}
  \vfill
  \begin{center}
    {\large\color{SecondaryColor}\itshape
      Let's talk about a property of fractals\ldots}

    \vspace{0.8cm}
    \begin{tikzpicture}
      \draw[PrimaryColor, line width=1.5pt] (-3.5,0) -- (3.5,0);
      \foreach \x in {-3,-2.5,...,3}
        \draw[AccentColor, line width=0.8pt] (\x,0.12) -- (\x,-0.12);
    \end{tikzpicture}

    \vspace{0.8cm}
    {\fontsize{26}{32}\selectfont\bfseries\color{PrimaryColor}
      Fractal Dimension}

    \vspace{0.5cm}
    {\small\color{DarkGray}
      What does it mean for a shape to have a
      \highlight{non-integer} dimension?}
  \end{center}
  \vfill
\end{frame}

% ---------------------------------------------------------------
% S2 Frame 2: Dimensions We Know (3 cumulative overlays)
% ---------------------------------------------------------------
\begin{frame}[shrink]{Dimensions We Know}
  \small
  We are all familiar with integer dimensions:
  \vspace{0.2cm}
  \begin{center}
    \begin{tikzpicture}[
        dbox/.style={rounded corners=3pt, draw=PrimaryColor,
          fill=PrimaryColor!12, inner sep=5pt,
          font=\Large\bfseries, text=PrimaryColor}
      ]
      % ── Row 1: Line → 1D ──
      \onslide<1->{
        \draw[ObjectColor, line width=3pt] (-4,2) -- (-1,2);
        \node[font=\footnotesize, text=DarkGray, below]
          at (-2.5,1.7) {Line};
        \draw[-{Stealth[length=5pt]}, thick, AccentColor]
          (-0.3,2) -- (0.5,2);
        \node[dbox] at (1.6,2) {1D};
      }

      % ── Row 2: Square → 2D ──
      \onslide<2->{
        \draw[ObjectColor, line width=1.5pt, fill=ObjectColor!18]
          (-4,-0.8) rectangle (-1.6,0.8);
        \node[font=\footnotesize, text=DarkGray, below]
          at (-2.8,-1) {Square};
        \draw[-{Stealth[length=5pt]}, thick, AccentColor]
          (-0.3,0) -- (0.5,0);
        \node[dbox] at (1.6,0) {2D};
      }

      % ── Row 3: Cube → 3D ──
      \onslide<3->{
        \begin{scope}[shift={(-4,-3.5)}]
          % front face
          \draw[ObjectColor, line width=1.5pt, fill=ObjectColor!20]
            (0,0) rectangle (1.8,1.8);
          % top face
          \draw[ObjectColor, line width=1.5pt, fill=ObjectColor!28]
            (0,1.8) -- (0.55,2.3) -- (2.35,2.3)
            -- (1.8,1.8) -- cycle;
          % right face
          \draw[ObjectColor, line width=1.5pt, fill=ObjectColor!12]
            (1.8,0) -- (2.35,0.5) -- (2.35,2.3)
            -- (1.8,1.8) -- cycle;
        \end{scope}
        \node[font=\footnotesize, text=DarkGray, below]
          at (-2.8,-3.75) {Cube};
        \draw[-{Stealth[length=5pt]}, thick, AccentColor]
          (-0.3,-2.6) -- (0.5,-2.6);
        \node[dbox] at (1.6,-2.6) {3D};
      }
    \end{tikzpicture}
  \end{center}

  \onslide<3->{
    \begin{center}
      {\itshape\color{DarkGray}
        Dimensions are always whole numbers\ldots\ right?}
    \end{center}
  }
\end{frame}

% ---------------------------------------------------------------
% S2 Frame 3: What About This One? (2 overlays)
% ---------------------------------------------------------------
\begin{frame}{What About This One?}
  \begin{columns}
    \begin{column}{0.55\textwidth}
      \only<1>{
        \small
        \vspace{0.5cm}
        {\large What dimension does this shape have?}

        \vspace{0.4cm}
        {\itshape Any guesses?}

        \vspace{0.3cm}
        \begin{itemize}
          \item It lives in 2D space\ldots
          \item But it has \highlight{zero area}
          \item And \highlight{infinite perimeter}
        \end{itemize}

        \vspace{0.3cm}
        It's not quite 2D, but it's clearly more than 1D\ldots
      }
      \only<2>{
        \begin{conceptbox}{Sierpinski Dimension}
          The Sierpinski Triangle has a fractal dimension of:
          \begin{center}
            {\fontsize{28}{34}\selectfont
              $D \approx \highlight{1.585}$}

            \vspace{0.1cm}
            {\Large\color{RayColor}\textbf{!!}}
          \end{center}

          \vspace{0.2cm}
          A \highlight{non-integer} dimension --- between a line
          (1D) and a filled surface (2D).
        \end{conceptbox}
      }
    \end{column}
    \begin{column}{0.4\textwidth}
      \begin{center}
        \begin{figure}
          \includegraphics[width=0.85\textwidth]
            {images/sierpinski/5.png}
          \only<1>{\caption*{Sierpinski Triangle}}
          \only<2>{\caption*{$D=\log3/\log2\approx1.585$}}
        \end{figure}
      \end{center}
    \end{column}
  \end{columns}
\end{frame}

% ---------------------------------------------------------------
% S2 Frame 4: Self-Similarity & Scaling (5 overlays)
% ---------------------------------------------------------------
\begin{frame}{Self-Similarity \& Scaling}
  \centering

  % ── Overlay 1: Scaling a Line ──
  \only<1>{%
    \begin{tcolorbox}[colback=PrimaryColor!8, colframe=PrimaryColor, arc=2mm, boxrule=1pt,
                      width=0.96\textwidth, halign=center, top=3pt, bottom=3pt]
      \textbf{\color{PrimaryColor}Scaling a Line ($D = 1$)}:\quad
      Scale by $r = \frac{1}{2} \implies$ \highlight{2 pieces}.
      \quad Scale again by $\frac{1}{2} \implies 2 \times 2 = \mathbf{4\text{ pieces}}$.
    \end{tcolorbox}

    \vspace{0.45cm}
    \begin{center}
      \begin{tikzpicture}[scale=0.88]
        % Original Line
        \draw[ObjectColor, line width=3.5pt] (-6.0, 0) -- (-3.4, 0);
        \node[font=\footnotesize\bfseries, PrimaryColor, below=4pt]
          at (-4.7, 0) {Original Line};

        % Arrow 1
        \draw[-{Stealth[scale=1.2]}, line width=2pt, AccentColor]
          (-2.9, 0) -- (-1.9, 0);
        \node[font=\small\bfseries, above=2pt, AccentColor]
          at (-2.4, 0.1) {$\times\frac{1}{2}$};

        % 2 Pieces
        \draw[ObjectColor, line width=3.5pt] (-1.4, 0) -- (-0.2, 0);
        \draw[ObjectColor, line width=3.5pt] (0.1, 0) -- (1.3, 0);
        \node[font=\footnotesize\bfseries, PrimaryColor, below=4pt]
          at (-0.05, 0) {\textbf{2} pieces};

        % Arrow 2
        \draw[-{Stealth[scale=1.2]}, line width=2pt, AccentColor]
          (1.8, 0) -- (2.8, 0);
        \node[font=\small\bfseries, above=2pt, AccentColor]
          at (2.3, 0.1) {$\times\frac{1}{2}$};

        % 4 Pieces
        \foreach \k in {0, 1, 2, 3} {
          \draw[ObjectColor, line width=3.5pt]
            (3.3 + \k*0.75, 0) -- ++(0.55, 0);
        }
        \node[font=\footnotesize\bfseries, PrimaryColor, below=4pt]
          at (4.7, 0) {\textbf{4} pieces};
      \end{tikzpicture}
    \end{center}
  }%

  % ── Overlay 2: Scaling a Square — Step 1 (4 pieces) ──
  \only<2>{%
    \begin{tcolorbox}[colback=PrimaryColor!8, colframe=PrimaryColor, arc=2mm, boxrule=1pt,
                      width=0.96\textwidth, halign=center, top=3pt, bottom=3pt]
      \textbf{\color{PrimaryColor}Scaling a Square ($D = 2$) --- Step 1}:\quad
      Scale by $r = \frac{1}{2} \implies$ \highlight{4 self-similar pieces} ($2 \times 2$ grid).
    \end{tcolorbox}

    \vspace{0.45cm}
    \begin{center}
      \begin{tikzpicture}[scale=0.85]
        % Original Square
        \draw[ObjectColor, line width=1.5pt, fill=ObjectColor!18]
          (-5.5, -0.9) rectangle ++(1.8, 1.8);
        \node[font=\footnotesize\bfseries, PrimaryColor, below=4pt]
          at (-4.6, -0.9) {Original Square};

        % Arrow 1
        \draw[-{Stealth[scale=1.1]}, line width=1.8pt, AccentColor]
          (-3.2, 0) -- (-2.2, 0);
        \node[font=\scriptsize\bfseries, above=1pt, AccentColor]
          at (-2.7, 0.1) {$\times\frac{1}{2}$};

        % 4 Pieces (2x2 grid)
        \foreach \i in {0, 1} {
          \foreach \j in {0, 1} {
            \draw[ObjectColor, line width=1pt, fill=ObjectColor!18]
              (-1.8 + \i*1.05, -0.85 + \j*1.05) rectangle ++(0.8, 0.8);
          }
        }
        \node[font=\footnotesize\bfseries, PrimaryColor, below=4pt]
          at (-0.95, -0.85) {\textbf{4} pieces};

        % Result Badge
        \node[font=\normalsize\bfseries, PrimaryColor, right]
          at (0.7, 0) {$N = 2^2 = \mathbf{4}$};
      \end{tikzpicture}
    \end{center}
  }%

  % ── Overlay 3: Scaling a Square — Step 2 (16 pieces) ──
  \only<3>{%
    \begin{tcolorbox}[colback=PrimaryColor!8, colframe=PrimaryColor, arc=2mm, boxrule=1pt,
                      width=0.96\textwidth, halign=center, top=3pt, bottom=3pt]
      \textbf{\color{PrimaryColor}Scaling a Square ($D = 2$) --- Step 2}:\quad
      Scale each piece by $\frac{1}{2}$ again $\implies 4 \times 4 = \mathbf{16\text{ self-similar pieces}}$ ($4 \times 4$ grid)!
    \end{tcolorbox}

    \vspace{0.45cm}
    \begin{center}
      \begin{tikzpicture}[scale=0.85]
        % Stage 1: Original Square
        \draw[ObjectColor, line width=1.5pt, fill=ObjectColor!18]
          (-5.5, -0.9) rectangle ++(1.8, 1.8);
        \node[font=\footnotesize\bfseries, PrimaryColor, below=4pt]
          at (-4.6, -0.9) {Original};

        % Arrow 1
        \draw[-{Stealth[scale=1.1]}, line width=1.8pt, AccentColor]
          (-3.2, 0) -- (-2.2, 0);
        \node[font=\scriptsize\bfseries, above=1pt, AccentColor]
          at (-2.7, 0.1) {$\times\frac{1}{2}$};

        % Stage 2: 4 Pieces (2x2 grid)
        \foreach \i in {0, 1} {
          \foreach \j in {0, 1} {
            \draw[ObjectColor, line width=1pt, fill=ObjectColor!18]
              (-1.8 + \i*1.05, -0.85 + \j*1.05) rectangle ++(0.8, 0.8);
          }
        }
        \node[font=\footnotesize\bfseries, PrimaryColor, below=4pt]
          at (-0.95, -0.85) {\textbf{4} pieces};

        % Arrow 2
        \draw[-{Stealth[scale=1.1]}, line width=1.8pt, AccentColor]
          (0.6, 0) -- (1.6, 0);
        \node[font=\scriptsize\bfseries, above=1pt, AccentColor]
          at (1.1, 0.1) {$\times\frac{1}{2}$};

        % Stage 3: 16 Pieces (4 clusters of 2x2 mini-squares)
        \foreach \ci in {0, 1} {
          \foreach \cj in {0, 1} {
            \foreach \i in {0, 1} {
              \foreach \j in {0, 1} {
                \draw[ObjectColor, line width=0.8pt, fill=ObjectColor!18]
                  (2.1 + \ci*1.02 + \i*0.44, -0.92 + \cj*1.02 + \j*0.44) rectangle ++(0.36, 0.36);
              }
            }
          }
        }
        \node[font=\footnotesize\bfseries, PrimaryColor, below=4pt]
          at (2.95, -0.92) {\textbf{16} pieces};

        % Result Badge
        \node[font=\normalsize\bfseries, PrimaryColor, right]
          at (4.5, 0) {$N = 4^2 = \mathbf{16}$};
      \end{tikzpicture}
    \end{center}
  }%

  % ── Overlay 4: Scaling a Cube ──
  \only<4>{%
    \begin{tcolorbox}[colback=PrimaryColor!8, colframe=PrimaryColor, arc=2mm, boxrule=1pt,
                      width=0.96\textwidth, halign=center, top=3pt, bottom=3pt]
      \textbf{\color{PrimaryColor}Scaling a Cube ($D = 3$)}:\quad
      Scale by $r = \frac{1}{2} \implies$ \highlight{8 self-similar pieces} ($2 \times 2 \times 2$ grid in 3D).
    \end{tcolorbox}

    \vspace{0.35cm}
    \begin{center}
      \begin{tikzpicture}[scale=0.80]
        % Left: Original Cube
        \filldraw[draw=ObjectColor, line width=1.5pt, fill=ObjectColor!30]
          (-4.2, 0.9) -- ++(0.65, 0.65) -- ++(1.8, 0) -- ++(-0.65, -0.65) -- cycle;
        \filldraw[draw=ObjectColor, line width=1.5pt, fill=ObjectColor!20]
          (-2.4, -0.9) -- ++(0.65, 0.65) -- ++(0, 1.8) -- ++(-0.65, -0.65) -- cycle;
        \filldraw[draw=ObjectColor, line width=1.5pt, fill=ObjectColor!14]
          (-4.2, -0.9) rectangle ++(1.8, 1.8);
        \node[font=\footnotesize\bfseries, PrimaryColor, below=4pt]
          at (-2.95, -1.1) {Original Cube};

        % Center: Transition Arrow
        \draw[-{Stealth[scale=1.2]}, line width=2pt, AccentColor]
          (-1.4, 0.1) -- (-0.1, 0.1);
        \node[font=\small\bfseries, above=2pt, AccentColor]
          at (-0.75, 0.25) {scale $\times\frac{1}{2}$};

        % Right: 8 Separated Small Cubes
        % Back layer (Z = 1, offset by (0.85, 0.85))
        \foreach \i in {0, 1} {
          \foreach \j in {0, 1} {
            \begin{scope}[shift={(0.6 + 0.85 + \i*1.30, -1.0 + 0.85 + \j*1.30)}]
              \filldraw[draw=ObjectColor, line width=1pt, fill=ObjectColor!30]
                (0, 0.75) -- ++(0.35, 0.35) -- ++(0.75, 0) -- ++(-0.35, -0.35) -- cycle;
              \filldraw[draw=ObjectColor, line width=1pt, fill=ObjectColor!20]
                (0.75, 0) -- ++(0.35, 0.35) -- ++(0, 0.75) -- ++(-0.35, -0.35) -- cycle;
              \filldraw[draw=ObjectColor, line width=1pt, fill=ObjectColor!14]
                (0, 0) rectangle ++(0.75, 0.75);
            \end{scope}
          }
        }

        % Front layer (Z = 0)
        \foreach \i in {0, 1} {
          \foreach \j in {0, 1} {
            \begin{scope}[shift={(0.6 + \i*1.30, -1.0 + \j*1.30)}]
              \filldraw[draw=ObjectColor, line width=1pt, fill=ObjectColor!30]
                (0, 0.75) -- ++(0.35, 0.35) -- ++(0.75, 0) -- ++(-0.35, -0.35) -- cycle;
              \filldraw[draw=ObjectColor, line width=1pt, fill=ObjectColor!20]
                (0.75, 0) -- ++(0.35, 0.35) -- ++(0, 0.75) -- ++(-0.35, -0.35) -- cycle;
              \filldraw[draw=ObjectColor, line width=1pt, fill=ObjectColor!14]
                (0, 0) rectangle ++(0.75, 0.75);
            \end{scope}
          }
        }

        \node[font=\footnotesize\bfseries, PrimaryColor, below=4pt]
          at (2.3, -1.1) {\textbf{8} separated cubes};

        % Result Badge
        \node[font=\normalsize\bfseries, PrimaryColor, right]
          at (4.3, 0.1) {$N = 2^3 = \mathbf{8}$};
      \end{tikzpicture}
    \end{center}
  }%

  % ── Overlay 5: Do You See a Pattern? ──
  \only<5>{%
    \begin{tcolorbox}[colback=PrimaryColor!8, colframe=PrimaryColor, arc=2mm, boxrule=1pt,
                      width=0.96\textwidth, halign=center, top=3pt, bottom=3pt]
      \textbf{\color{PrimaryColor}Do You See a Pattern?}\\
      Each scaling step multiplies pieces by $(1/r)^D$!
    \end{tcolorbox}

    \vspace{0.25cm}
    \begin{columns}
      \begin{column}{0.55\textwidth}
        \centering
        \small
        \renewcommand{\arraystretch}{1.3}
        \begin{tabular}{lcccc}
          \toprule
          \textbf{Shape} & \textbf{Scale} & \textbf{1st Step} & \textbf{2nd Step} & \textbf{$D$} \\
          \midrule
          Line   & $1/2$ & $2 = 2^1$ & $4 = 2^2$  & \textbf{1} \\
          Square & $1/2$ & $4 = 2^2$ & $16 = 2^4$ & \textbf{2} \\
          Cube   & $1/2$ & $8 = 2^3$ & $(64 = 2^6)$ & \textbf{3} \\
          \bottomrule
        \end{tabular}

        \vspace{0.3cm}
        \begin{conceptbox}{The Pattern}
          \centering
          \highlight{Number of pieces} $= \left(\dfrac{1}{r}\right)^{\!D}$
        \end{conceptbox}
      \end{column}
      \begin{column}{0.4\textwidth}
        \begin{center}
          \includegraphics[width=0.45\textwidth]{images/sierpinski/3.png}

          \vspace{0.2cm}
          {\Large \textbf{\color{PrimaryColor}$D = \;?$}}\\
          \vspace{0.1cm}
          \footnotesize What if the shape is a \highlight{fractal}?
        \end{center}
      \end{column}
    \end{columns}
  }%
\end{frame}

% ---------------------------------------------------------------
% S2 Frame 5: The Self-Similarity Rule (2 overlays)
% ---------------------------------------------------------------
\begin{frame}[shrink]{The Self-Similarity Rule}
  \centering
  \small
  After scaling by factor $r$, the number of self-similar pieces is:
  \quad
  \highlight{\Large $N = \left(\dfrac{1}{r}\right)^{\!D}$}

  \vspace{0.25cm}
  \only<1>{%
    {\normalsize \textbf{\color{PrimaryColor}Square ($D = 2$)} \quad with scale factor $r = \frac{1}{2}$:\quad
    $N = \left(\frac{1}{1/2}\right)^{\!2} = 2^2 = \mathbf{4\text{ pieces}}$}%
  }%
  \only<2>{%
    {\normalsize \textbf{\color{PrimaryColor}Cube ($D = 3$)} \quad with scale factor $r = \frac{1}{2}$:\quad
    $N = \left(\frac{1}{1/2}\right)^{\!3} = 2^3 = \mathbf{8\text{ pieces}}$}%
  }%

  \vspace{0.3cm}
  \only<1>{%
    \begin{center}
      \begin{tikzpicture}[scale=0.85]
        % ── Left: Original Square ──
        \draw[ObjectColor, line width=1.5pt, fill=ObjectColor!18]
          (-4.2, -1.0) rectangle ++(2.0, 2.0);
        \node[font=\footnotesize\bfseries, PrimaryColor, below]
          at (-3.2, -1.25) {Original Square};

        % ── Center: Transition Arrow ──
        \draw[-{Stealth[scale=1.2]}, line width=2pt, AccentColor]
          (-1.4, 0) -- (-0.1, 0);
        \node[font=\small\bfseries, above=2pt, AccentColor]
          at (-0.75, 0.15) {scale $\times\frac{1}{2}$};

        % ── Right: 4 Separated Squares ──
        \foreach \i in {0, 1} {
          \foreach \j in {0, 1} {
            \draw[ObjectColor, line width=1.2pt, fill=ObjectColor!18]
              (0.8 + \i*1.25, -1.075 + \j*1.25) rectangle ++(0.9, 0.9);
          }
        }
        \node[font=\footnotesize\bfseries, PrimaryColor, below]
          at (1.875, -1.25) {\textbf{4} separated pieces};

        % ── Result Badge ──
        \node[font=\normalsize\bfseries, PrimaryColor, right]
          at (3.6, 0) {$N = 2^2 = \mathbf{4}$};
      \end{tikzpicture}
    \end{center}
  }%
  \only<2>{%
    \begin{center}
      \begin{tikzpicture}[scale=0.80]
        % ── Left: Original Cube ──
        \filldraw[draw=ObjectColor, line width=1.5pt, fill=ObjectColor!30]
          (-4.2, 0.9) -- ++(0.65, 0.65) -- ++(1.8, 0) -- ++(-0.65, -0.65) -- cycle;
        \filldraw[draw=ObjectColor, line width=1.5pt, fill=ObjectColor!20]
          (-2.4, -0.9) -- ++(0.65, 0.65) -- ++(0, 1.8) -- ++(-0.65, -0.65) -- cycle;
        \filldraw[draw=ObjectColor, line width=1.5pt, fill=ObjectColor!14]
          (-4.2, -0.9) rectangle ++(1.8, 1.8);
        \node[font=\footnotesize\bfseries, PrimaryColor, below]
          at (-2.95, -1.25) {Original Cube};

        % ── Center: Transition Arrow ──
        \draw[-{Stealth[scale=1.2]}, line width=2pt, AccentColor]
          (-1.4, 0.1) -- (-0.1, 0.1);
        \node[font=\small\bfseries, above=2pt, AccentColor]
          at (-0.75, 0.25) {scale $\times\frac{1}{2}$};

        % ── Right: 8 Separated Small Cubes ──
        % Back layer (Z = 1, offset by (0.85, 0.85))
        \foreach \i in {0, 1} {
          \foreach \j in {0, 1} {
            \begin{scope}[shift={(0.6 + 0.85 + \i*1.30, -1.0 + 0.85 + \j*1.30)}]
              \filldraw[draw=ObjectColor, line width=1pt, fill=ObjectColor!30]
                (0, 0.75) -- ++(0.35, 0.35) -- ++(0.75, 0) -- ++(-0.35, -0.35) -- cycle;
              \filldraw[draw=ObjectColor, line width=1pt, fill=ObjectColor!20]
                (0.75, 0) -- ++(0.35, 0.35) -- ++(0, 0.75) -- ++(-0.35, -0.35) -- cycle;
              \filldraw[draw=ObjectColor, line width=1pt, fill=ObjectColor!14]
                (0, 0) rectangle ++(0.75, 0.75);
            \end{scope}
          }
        }

        % Front layer (Z = 0)
        \foreach \i in {0, 1} {
          \foreach \j in {0, 1} {
            \begin{scope}[shift={(0.6 + \i*1.30, -1.0 + \j*1.30)}]
              \filldraw[draw=ObjectColor, line width=1pt, fill=ObjectColor!30]
                (0, 0.75) -- ++(0.35, 0.35) -- ++(0.75, 0) -- ++(-0.35, -0.35) -- cycle;
              \filldraw[draw=ObjectColor, line width=1pt, fill=ObjectColor!20]
                (0.75, 0) -- ++(0.35, 0.35) -- ++(0, 0.75) -- ++(-0.35, -0.35) -- cycle;
              \filldraw[draw=ObjectColor, line width=1pt, fill=ObjectColor!14]
                (0, 0) rectangle ++(0.75, 0.75);
            \end{scope}
          }
        }

        \node[font=\footnotesize\bfseries, PrimaryColor, below]
          at (2.3, -1.25) {\textbf{8} separated small cubes};

        % ── Result Badge ──
        \node[font=\normalsize\bfseries, PrimaryColor, right]
          at (4.3, 0.1) {$N = 2^3 = \mathbf{8}$};
      \end{tikzpicture}
    \end{center}

    \vspace{-0.05cm}
    \centering\footnotesize
    \textit{Can we use this same rule to find the dimension of a \highlight{fractal}?}
  }%
\end{frame}

% ---------------------------------------------------------------
% S2 Frame 6: Dimension of the Sierpinski Triangle (4 overlays)
% ---------------------------------------------------------------
\begin{frame}{Dimension of the Sierpinski Triangle}
  \begin{columns}
    \begin{column}{0.55\textwidth}
      \small
      \only<1>{
        Let's apply the self-similarity rule to the
        \highlight{Sierpinski Triangle}.

        \vspace{0.3cm}
        \begin{conceptbox}{Step 1}
          Scale the Sierpinski Triangle by $\frac{1}{2}$.
        \end{conceptbox}
      }
      \only<2>{
        \begin{conceptbox}{Step 2 --- Self-Similar Pieces}
          We get exactly \highlight{3} self-similar copies
          of the original triangle.

          \vspace{0.2cm}
          Each smaller triangle is a perfect replica,
          scaled by $r = \frac{1}{2}$.
        \end{conceptbox}
      }
      \only<3>{
        \begin{mathbox}{Step 3 --- The Equation}
          Using $N = \left(\frac{1}{r}\right)^D$:
          \begin{align*}
            3 &= \left(\frac{1}{1/2}\right)^{\!D} \\[6pt]
            3 &= 2^D
          \end{align*}
        \end{mathbox}
      }
      \only<4>{
        \begin{mathbox}{Step 4 --- Solving for $D$}
          \begin{align*}
            2^D &= 3 \\[4pt]
            D\,\log 2 &= \log 3 \\[4pt]
            D &= \frac{\log 3}{\log 2}
          \end{align*}
          \begin{center}
            {\fontsize{22}{28}\selectfont
              $D \approx \highlight{1.585}$}
          \end{center}
        \end{mathbox}

        \vspace{0.1cm}
        \highlight{Not an integer!} --- strictly between
        1D (line) and 2D (surface).
      }
    \end{column}
    \begin{column}{0.4\textwidth}
      \begin{center}
        \begin{tikzpicture}[scale=0.85]
          \only<1>{
            \node at (0,0) {
              \includegraphics[width=0.85\textwidth]
                {images/sierpinski/5.png}
            };
            \draw[-{Stealth}, very thick, AccentColor]
              (0,-1.8) -- (0,-2.5);
            \node[font=\small, AccentColor, right]
              at (0.1,-2.15) {scale by $\frac{1}{2}$};
          }
          \only<2-3>{
            % Three self-similar copies
            \node at (0,1.2) {
              \includegraphics[width=0.4\textwidth]
                {images/sierpinski/4.png}
            };
            \node at (-1,-0.6) {
              \includegraphics[width=0.4\textwidth]
                {images/sierpinski/4.png}
            };
            \node at (1,-0.6) {
              \includegraphics[width=0.4\textwidth]
                {images/sierpinski/4.png}
            };
            \node[font=\Large\bfseries, PrimaryColor]
              at (0,-2) {3 copies};
          }
          \only<4>{
            \node at (0,0.5) {
              \includegraphics[width=0.7\textwidth]
                {images/sierpinski/5.png}
            };
            \node[rounded corners=4pt, draw=PrimaryColor,
              fill=PrimaryColor!12, inner sep=6pt,
              font=\Large\bfseries, text=PrimaryColor]
              at (0,-2) {$D\approx1.585$};
          }
        \end{tikzpicture}
      \end{center}
    \end{column}
  \end{columns}
\end{frame}

% ---------------------------------------------------------------
% S2 Frame 7: Similarly — Koch Curve (2 overlays)
% ---------------------------------------------------------------
\begin{frame}{Similarly: The Koch Curve}
  \begin{columns}
    \begin{column}{0.55\textwidth}
      \small
      \only<1>{
        The \highlight{Koch Curve} can be divided into
        \highlight{4} self-similar pieces, each scaled by
        $r = \frac{1}{3}$.

        \vspace{0.3cm}
        \begin{conceptbox}{Setup}
          \begin{itemize}
            \item Scale factor: $r = \frac{1}{3}$
            \item Self-similar pieces: $N = 4$
          \end{itemize}
        \end{conceptbox}
      }
      \only<2>{
        \begin{mathbox}{Koch Dimension}
          \begin{align*}
            N &= \left(\frac{1}{r}\right)^{\!D} \\[4pt]
            4 &= 3^D \\[4pt]
            D &= \frac{\log 4}{\log 3}
          \end{align*}
          \begin{center}
            {\fontsize{22}{28}\selectfont
              $D \approx \highlight{1.262}$}
          \end{center}
        \end{mathbox}

        \vspace{0.2cm}
        Also \highlight{non-integer} --- between a line (1D)
        and a surface (2D), but closer to a line than
        the Sierpinski Triangle.
      }
    \end{column}
    \begin{column}{0.4\textwidth}
      \begin{center}
        \begin{tikzpicture}
          \only<1>{
            \node at (0,1.5) {
              \includegraphics[width=0.85\textwidth]
                {images/koch/5.png}
            };
            \draw[-{Stealth}, very thick, AccentColor]
              (0,0.3) -- (0,-0.4);
            \node[font=\small, AccentColor, right]
              at (0.1,-0.05) {$\times\frac{1}{3}$};
            % 4 smaller copies arranged in Koch shape
            \node at (-1.5,-1.5) {
              \includegraphics[width=0.22\textwidth]
                {images/koch/4.png}
            };
            \node[rotate=60] at (-0.45,-1.2) {
              \includegraphics[width=0.22\textwidth]
                {images/koch/4.png}
            };
            \node[rotate=-60] at (0.45,-1.2) {
              \includegraphics[width=0.22\textwidth]
                {images/koch/4.png}
            };
            \node at (1.5,-1.5) {
              \includegraphics[width=0.22\textwidth]
                {images/koch/4.png}
            };
            \node[font=\Large\bfseries, PrimaryColor]
              at (0,-2.8) {4 copies};
          }
          \only<2>{
            \node at (0,0.5) {
              \includegraphics[width=0.85\textwidth]
                {images/koch/5.png}
            };
            \node[rounded corners=4pt, draw=PrimaryColor,
              fill=PrimaryColor!12, inner sep=6pt,
              font=\Large\bfseries, text=PrimaryColor]
              at (0,-1.5) {$D\approx1.262$};
          }
        \end{tikzpicture}
      \end{center}
    \end{column}
  \end{columns}
\end{frame}

% ---------------------------------------------------------------
% S2 Frame 8: The General Theory — Box-Counting Dimension
% ---------------------------------------------------------------
\begin{frame}[shrink]{The General Theory: Box-Counting Dimension}
  \small
  The \highlight{Self-Similarity Dimension} works beautifully for
  self-similar fractals. But what about more complex shapes?

  \vspace{0.3cm}
  \begin{mathbox}{Box-Counting Dimension}
    The fractal dimension $D$ is defined as:
    $$D = \lim_{\varepsilon \to 0}
      \frac{\log N(\varepsilon)}{\log(1/\varepsilon)}$$
    where $N(\varepsilon)$ is the number of boxes of side length
    $\varepsilon$ needed to cover the fractal.
  \end{mathbox}

  \vspace{0.3cm}
  \pause
  \begin{conceptbox}{Intuition}
    \begin{itemize}
      \item Cover the fractal with a grid of boxes of size
        $\varepsilon$
      \item Count how many boxes $N(\varepsilon)$ the fractal
        touches
      \item Shrink $\varepsilon \to 0$ and observe how
        $N(\varepsilon)$ grows
      \item The \highlight{rate of growth} tells us the
        dimension
    \end{itemize}
  \end{conceptbox}

  \vspace{0.2cm}
  \begin{center}
    {\footnotesize\color{DarkGray} For a deeper dive, see
    \href{https://www.youtube.com/watch?v=gB9n2gHsHN4}%
    {\textcolor{AccentColor}{3Blue1Brown:
    Fractals are typically not self-similar}}}
  \end{center}
\end{frame}

% ---------------------------------------------------------------
% S2 Frame 9: Handoff to Speaker 3
% ---------------------------------------------------------------
\begin{frame}[shrink]{Handoff to Speaker 3}
  \small
  We've now uncovered a remarkable property of fractals:

  \begin{itemize}
    \item They have \highlight{non-integer dimensions} ---
      genuinely ``between'' the familiar 1D, 2D, 3D.
    \item The Sierpinski Triangle lives in dimension
      $D\approx 1.585$.
    \item The Koch Curve lives in dimension
      $D\approx 1.262$.
    \item The \highlight{box-counting dimension} generalises
      this to \emph{any} shape.
  \end{itemize}

  \pause
  \vspace{0.3cm}
  \begin{center}
    \begin{tikzpicture}
      \node[conceptbox, text width=0.8\textwidth] {
        \centering
        We now know \emph{what} fractals are and how to
        measure their complexity. \\[4pt]
        But \highlight{where do we see them}? And
        \highlight{how are they used} in computer science?
      };
    \end{tikzpicture}
  \end{center}

  \vspace{0.05cm}
  \pause
  That's exactly where Speaker~3 takes us ---
  fractals in nature, art, and computation.
\end{frame}

% ================================================================
% PART 3 — FRACTALS IN THE REAL WORLD (Speaker 3)
% 9 frames: section title, nature, coastline (2), body, database,
%   game worlds, limits, conclusion  (+ references)
% ================================================================
\section{Fractals in the Real World}

% ---------------------------------------------------------------
% S3 Frame 1: Transition — section title (same style as Speaker 2)
% ---------------------------------------------------------------
\begin{frame}{}
  \vfill
  \begin{center}
    {\large\color{SecondaryColor}\itshape
      We know what a fractal is, and how to measure it\ldots}

    \vspace{0.8cm}
    \begin{tikzpicture}
      \draw[PrimaryColor, line width=1.5pt] (-3.5,0) -- (3.5,0);
      \foreach \x in {-3,-2.5,...,3}
        \draw[AccentColor, line width=0.8pt] (\x,0.12) -- (\x,-0.12);
    \end{tikzpicture}

    \vspace{0.8cm}
    {\fontsize{26}{32}\selectfont\bfseries\color{PrimaryColor}
      Fractals in the Real World}

    \vspace{0.5cm}
    {\small\color{DarkGray}
      Where does all of this actually \highlight{show up}?}
  \end{center}
  \vfill
\end{frame}

% ---------------------------------------------------------------
% S3 Frame 2: Nature — real photos next to a rule-grown tree
% ---------------------------------------------------------------
\begin{frame}{Nature draws with one rule, repeated}
  \centering
  \vspace{0.04cm}
  \includegraphics[width=0.285\textwidth]{images/part3/nature_romanesco.png}%
  \hspace{0.05\textwidth}%
  \onslide<2->{\includegraphics[width=0.285\textwidth]{images/part3/nature_snowflake.png}}\\[7pt]
  \onslide<3->{\includegraphics[width=0.285\textwidth]{images/part3/nature_aloe.png}}%
  \hspace{0.05\textwidth}%
  \onslide<4->{\includegraphics[width=0.285\textwidth]{images/part3/nature_pinecone.png}}

  \vspace{0.18cm}
  \onslide<5->{\takeaway{Broccoli, ice, leaves, cones --- the \highlight{same shape at every zoom}.}}
\end{frame}

% ---------------------------------------------------------------
% S3 Frame 3: Pythagorean tree — a plant grown from one square rule
% ---------------------------------------------------------------
\begin{frame}{A tree built from nothing but squares}
  \centering
  \includegraphics[width=0.31\textwidth]{images/part3/pythagoras_1.png}\hfill
  \onslide<2->{\includegraphics[width=0.31\textwidth]{images/part3/pythagoras_2.png}}\hfill
  \onslide<3->{\includegraphics[width=0.31\textwidth]{images/part3/pythagoras_3.png}}

  \vspace{0.2cm}
  \onslide<4->{{\small Stand a \highlight{right triangle} on each square's top edge and grow a
    new square on each leg --- so $a^2+b^2=c^2$ sits at every single joint.}}

  \vspace{0.2cm}
  \onslide<5->{\takeaway{Nobody designed a single branch. \highlight{One rule}, repeated twelve
    times, and a tree appears.}}

  \source{Pythagorean tree, after A.~E.~Bosman (1942), drawn at $45^\circ$.}
\end{frame}

% ---------------------------------------------------------------
% S3 Frame 3b: the same rule, tilted — now it looks alive
% ---------------------------------------------------------------
\begin{frame}{Tilt the rule and it starts to look alive}
  \begin{columns}[c]
    \begin{column}{0.57\textwidth}
      \centering
      \includegraphics[width=\textwidth]{images/part3/pythagoras_4.png}
    \end{column}
    \begin{column}{0.41\textwidth}
      \small
      \begin{itemize}
        \item<2-> At $45^\circ$ the two branches match \highlight{exactly} --- too perfect to be real
        \item<3-> Tilt that triangle to $55^\circ$ and they stop matching
        \item<4-> The tree \highlight{leans}, the way a real one grows toward the light
      \end{itemize}
      \vspace{0.3cm}
      \onslide<5->{\takeaway{\highlight{One number changed} --- and a diagram turns into a plant.}}
    \end{column}
  \end{columns}
\end{frame}

% ---------------------------------------------------------------
% S3 Frame 5: The body — branching packs area into volume
% ---------------------------------------------------------------
\begin{frame}{Branching packs a huge surface into a small space}
  \centering
  \includegraphics[height=0.25\textwidth]{images/part3/lungs.png}\hfill
  \onslide<2->{\includegraphics[height=0.25\textwidth]{images/part3/eye_full.png}}\hfill
  \onslide<3->{\includegraphics[height=0.25\textwidth]{images/part3/eye_zoom.png}}

  \vspace{0.12cm}
  {\scriptsize\color{DarkGray} airways and blood vessels of the lung \quad
   vessels in the retina \quad the same vessels, zoomed in}

  \vspace{0.15cm}
  \begin{itemize}\small
    \item<4-> Lungs: about \highlight{23 splits} $\Rightarrow$ $\sim$\highlight{70 m$^2$}
      of surface packed inside $\sim$5 litres
    \item<5-> Doctors track the \highlight{dimension} of retinal vessels to spot disease early
  \end{itemize}

  \vspace{0.05cm}
  \onslide<6->{\takeaway{Branching again and again is nature's cheapest way to get \highlight{area out of volume}.}}

  \source{Lung branching figures: Weibel (1963).}
\end{frame}

% ---------------------------------------------------------------
% S3 Frame 6b: Fractal compression — store the rule, not the pixels
% ---------------------------------------------------------------
\begin{frame}{Fractal compression: a whole plant in 28 numbers}
  \begin{columns}[c]
    \begin{column}{0.4\textwidth}
      \centering
      \only<1>{\includegraphics[width=0.85\textwidth]{images/part3/fern_1.png}}%
      \only<2>{\includegraphics[width=0.85\textwidth]{images/part3/fern_2.png}}%
      \only<3->{\includegraphics[width=0.85\textwidth]{images/part3/fern_3.png}}
    \end{column}
    \begin{column}{0.58\textwidth}
      \small
      \begin{itemize}
        \item<2-> Four rules, each a \highlight{shrink-and-tilt} of the whole picture, applied at random
        \item<3-> $4\times 6$ coefficients $+$ 4 probabilities $=$ \highlight{28 numbers}
        \item<4-> Those 28 numbers redraw the fern at \highlight{any resolution} --- zoom in, there is always more leaf
      \end{itemize}
      \vspace{0.2cm}
      \onslide<5->{\takeaway{Store the \highlight{rule}, not the pixels --- the idea behind
        \highlight{fractal image compression}.}}
    \end{column}
  \end{columns}

  \source{Barnsley fern: an iterated function system (Barnsley, \textit{Fractals Everywhere}, 1988), plotted by the chaos game.}
\end{frame}

% ---------------------------------------------------------------
% S3 Frame 7: Game worlds — grown from a rule, not stored
% ---------------------------------------------------------------
\begin{frame}{Game worlds are grown from a rule, not stored}
  \centering
  \includegraphics[width=0.32\textwidth]{images/part3/terrain_octaves_1.png}\hfill
  \onslide<2->{\includegraphics[width=0.32\textwidth]{images/part3/terrain_octaves_2.png}}\hfill
  \onslide<3->{\includegraphics[width=0.32\textwidth]{images/part3/terrain_octaves_3.png}}

  \vspace{0.15cm}
  \onslide<4->{\takeaway{Add \highlight{smaller bumps at every scale} --- the Koch idea --- and a
    handful of numbers becomes an endless mountain range. Store the rule, not the map.}}

  \source{Terrain generated here by summing 7 layers of noise, each half the size of the last (fractal Brownian motion).}
\end{frame}

% ---------------------------------------------------------------
% S3 Frame 7b: how far the same rule can go
% ---------------------------------------------------------------
\begin{frame}{Keep stacking layers and you get a world}
  \centering
  \includegraphics[width=0.82\textwidth]{images/part3/voxel_world.png}

  \vspace{0.05cm}
  \onslide<2->{\takeaway{Round the same heights off to whole \highlight{blocks} and you get a
    sandbox world --- grass, beaches, lakes, a snow peak. Still just the rule.}}

  \source{The same 15 layers of noise, rounded into blocks and drawn isometrically.}
\end{frame}

% ---------------------------------------------------------------
% S3 Frame 8: Limits — real fractals stop
% ---------------------------------------------------------------
\begin{frame}{One honest catch: real fractals stop}
  \vspace{0.3cm}
  \large
  \begin{itemize}\itemsep0.45em
    \item<1-> A fern looks like a fern \ldots\ until you reach a single \highlight{cell}
    \item<2-> Lungs stop branching after about \highlight{23 splits}
    \item<3-> A tree runs out at the last \highlight{twig}
    \item<4-> Only the mathematical ones --- Sierpinski, Koch --- go on \highlight{forever}
  \end{itemize}

  \vspace{0.5cm}
  \normalsize
  \onslide<5->{\takeaway{A fractal is a \highlight{model} of nature --- and every good model knows
    the range where it works.}}
\end{frame}

% ---------------------------------------------------------------
% References
% ---------------------------------------------------------------
\begin{frame}{References}
  \footnotesize
  \begin{itemize}
    \item B.~Mandelbrot, ``How long is the coast of Britain? Statistical self-similarity and fractional dimension,'' \textit{Science}, 1967.
    \item B.~Mandelbrot, \textit{The Fractal Geometry of Nature}, 1982.
    \item E.~R.~Weibel, \textit{Morphometry of the Human Lung}, Springer, 1963.
    \item M.~F.~Barnsley, \textit{Fractals Everywhere}, Academic Press, 1988 (iterated function systems, the fern).
    \item A.~E.~Bosman, \textit{Het wondere onderzoekingsveld van de vlakke meetkunde}, 1957 (Pythagorean tree, drawn 1942).
    \item A.~Fournier, D.~Fussell, L.~Carpenter, ``Computer rendering of stochastic models,'' \textit{Comm.~ACM}, 1982 (fractal terrain).
    \item 3Blue1Brown, ``Fractals are typically not self-similar'' (video).
  \end{itemize}
\end{frame}

\end{document}
